\chapter*{Curriculum connections and sources}
\addcontentsline{toc}{chapter}{Curriculum connections and sources}
The main learning sequence is designed as enrichment for senior students in NSW Mathematics Extension, VCE Specialist Mathematics and QCE Specialist Mathematics. These courses offer useful preparation in proof, algebra, counting or matrices, but their emphases differ; the booklet therefore introduces graph theory directly and includes a matrix primer. The applied appendix gathers network ideas connected to NSW Standard, QCE General and VCE General Mathematics.

The course tags in Appendix 1 refer to the specific editions below. The mapping was checked against these sources on 16 October 2026; syllabus versions and implementation dates can change, so consult the relevant authority for current requirements. This booklet is an enrichment resource, not an official syllabus or assessment guide.

\section*{Curriculum sources}
\begin{itemize}
\item \textbf{NSW:} NSW Education Standards Authority (NESA), \emph{Mathematics Standard 11--12 Syllabus (2024)}. [NSW11] identifies the Year 11 ``Networks, paths and trees'' content for Standard 1 and Standard 2; [NSW12] identifies Year 12 Standard 2 content, including critical path analysis and network flow. See the \href{https://curriculum.nsw.edu.au/learning-areas/mathematics/mathematics-standard-11-12-2024/overview/course}{official course overview} and \href{https://curriculum.nsw.edu.au/learning-areas/mathematics/mathematics-standard-11-12-2024/content/year-12/fa493dfdfb}{Year 12 content page}.
\item \textbf{Queensland:} Queensland Curriculum and Assessment Authority (QCAA), \emph{General Mathematics 2025 v1.3 General senior syllabus}. [QCE] denotes Unit 4. See the \href{https://www.qcaa.qld.edu.au/downloads/senior-qce/syllabuses/snr_maths_general_25_syll.pdf}{official syllabus}.
\item \textbf{Victoria:} Victorian Curriculum and Assessment Authority (VCAA), \emph{VCE General Mathematics Study Design 2023}. [VCE2] and [VCE4] denote Units 2 and 4, respectively. See the VCAA \href{https://vcaa.vic.edu.au/curriculum/vce-curriculum/vce-study-designs/general-mathematics/vce-general-mathematics}{study-design page} for the official 2023 Mathematics Study Design, and its \href{https://www.vcaa.vic.edu.au/curriculum/vce-curriculum/vce-study-designs/general-mathematics/planning}{planning page} for the course structure and topic sequence.
\end{itemize}

\section*{Source for the four-colour theorem}
The historical statement in the graph-colouring chapter is based on the original announcement and the detailed proof papers:
\begin{itemize}
\item K. Appel and W. Haken, ``Every planar map is four colorable,'' \emph{Bulletin of the American Mathematical Society} 82, no. 5 (1976), 711--712. \href{https://doi.org/10.1090/S0002-9904-1976-14122-5}{doi:10.1090/S0002-9904-1976-14122-5}.
\item K. Appel and W. Haken, ``Every planar map is four colorable. Part I. Discharging,'' \emph{Illinois Journal of Mathematics} 21, no. 3 (1977), 429--490. \href{https://doi.org/10.1215/ijm/1256049011}{doi:10.1215/ijm/1256049011}.
\item K. Appel, W. Haken, and J. Koch, ``Every planar map is four colorable. Part II. Reducibility,'' \emph{Illinois Journal of Mathematics} 21, no. 3 (1977), 491--567. \href{https://doi.org/10.1215/ijm/1256049012}{doi:10.1215/ijm/1256049012}.
\end{itemize}
